\documentclass[tikz,border=6pt]{standalone}
\usepackage{tikz}
\usetikzlibrary{positioning,arrows.meta}
\usepackage{graphicx}
\usepackage{xcolor}

\definecolor{gitopsblue}{RGB}{51,108,229}
\definecolor{imperativegray}{RGB}{102,102,102}
\definecolor{metricyellow}{RGB}{251,192,45}

\begin{document}
\begin{tikzpicture}[
    node distance=0.30cm,
    box/.style={rectangle, rounded corners=3pt, draw=#1, fill=#1!8,
                text width=2.45cm, minimum height=2.0cm, align=center,
                line width=0.8pt, font=\scriptsize, inner sep=3pt},
    note/.style={rectangle, rounded corners=3pt, draw=gray!45, fill=gray!4,
                 text width=4.0cm, align=left, font=\scriptsize, inner sep=5pt},
    arrow/.style={-{Latex[length=2mm]}, line width=0.9pt, draw=gray!65}
]

\node[box=gitopsblue, text width=2.5cm] (stack) {
    \includegraphics[height=9pt]{logos/kubernetes.png}\\[1pt]
    \textbf{Satu \textit{stack}}\\
    InvenioRDM\\
    27 objek render
};
\node[box=gitopsblue, right=of stack, text width=2.7cm] (gitops) {
    \includegraphics[height=9pt]{logos/argocd.png}\\[1pt]
    \textbf{Mode GitOps}\\
    Amati keadaan aktual\\
    ArgoCD aktif\\
    7$\times$15 percobaan
};
\node[box=imperativegray, right=of gitops, text width=2.45cm] (reset) {
    \textbf{Pemulihan}\\
    \textbf{kondisi awal}\\
    verifikasi inventaris\\
    jeda 60 detik
};
\node[box=imperativegray, right=of reset, text width=2.9cm] (imperative) {
    \includegraphics[height=9pt]{logos/kubernetes.png}\\[1pt]
    \textbf{Mode imperatif}\\
    Amati keadaan aktual\\
    ArgoCD ditangguhkan\\
    pemeriksa 5 detik\\
    + \textit{runbook}
};
\node[box=metricyellow!85!black, right=of imperative, text width=2.9cm] (eval) {
    \textbf{Evaluator}\\
    deteksi + pemulihan\\
    kecocokan konfigurasi\\
    kesehatan: pemeriksaan pengaman
};

\draw[arrow] (stack) -- (gitops);
\draw[arrow] (gitops) -- (reset);
\draw[arrow] (reset) -- (imperative);
\draw[arrow] (imperative) -- (eval);

\node[note, below=0.42cm of stack.south west, anchor=north west] (n1)
    {\textbf{Bukan dua \textit{deployment}.}\\Satu \textit{stack} dipakai bergantian.};
\node[note, below=0.42cm of reset.south, anchor=north] (n2)
    {\textbf{HPA aktif pada kedua mode.}\\Skenario A adalah kontrol HPA.};
\node[note, below=0.42cm of eval.south east, anchor=north east] (n3)
    {\textbf{Tanpa waktu reaksi manusia.}\\Urutan mode tetap dan tidak diacak.};

\end{tikzpicture}
\end{document}
